% TOP_PROBLEM: {"id": 2800301, "title": "10 Lectures and 42 Open Problems — Hardness at the Cheeger square-root gap", "category_id": 15, "category": {"id": 15, "name": "computer_science", "display_name": "Computer Science", "description": "Computational complexity, algorithms, and theoretical CS.", "slug": "computer-science", "order_index": 15, "created_at": "2026-07-31T15:26:25.671Z"}, "status": "partially_solved", "problem_number": "AMR-027-0301", "set_id": 13}
% FIRST_POSTED: 2026-10-02
% ENTRY_KIND: statement_only
% UnsolvedMath ID 2800301; source catalogue code AMR-027-0301
% !TeX program = lualatex
% Definition and sources only; this file is not an LLM solution attempt.
\documentclass[11pt,a4paper]{article}
\usepackage[margin=25mm]{geometry}
\usepackage{fontspec}
\setmainfont{lmroman10-regular.otf}[BoldFont=lmroman10-bold.otf,
 ItalicFont=lmroman10-italic.otf,BoldItalicFont=lmroman10-bolditalic.otf]
\usepackage{amsmath,amssymb,mathtools}
\usepackage{unicode-math}
\setmathfont{latinmodern-math.otf}
\AtBeginDocument{\let\setminus\smallsetminus}
\usepackage{xurl,hyperref}
\hypersetup{colorlinks=true,urlcolor=blue}
\setcounter{secnumdepth}{0}
\setlength{\parindent}{0pt}
\setlength{\parskip}{5pt}
\setlength{\emergencystretch}{3em}
\begin{document}
\section*{10 Lectures and 42 Open Problems — Hardness at the Cheeger square-root gap}\label{2800301}
{\small UnsolvedMath ID 2800301 · AMR-027-0301 · Computer Science · AMR Open Problem Lists}
\subsection{Definitions and mathematical statement}
Let $G=(V,E)$ be a finite undirected graph without isolated vertices.
Write $d(v)$ for its degree, $\operatorname{vol}(S)=\sum_{v\in S}d(v)$,
and $\delta(S)$ for the edges with one endpoint in $S$ and one outside.
Its conductance, or normalized Cheeger constant, is
\[
 h_G=\min_{0<\operatorname{vol}(S)\le\operatorname{vol}(V)/2}
                   \frac{|\delta(S)|}{\operatorname{vol}(S)}.
\]
The minimization ranges over all eligible vertex sets, without a
prescribed cardinality or a requirement that the two sides be equal.

Does there exist an absolute constant $c>0$ for which the following
promise problem is NP-hard? The input is $G$ and a positive rational
$\phi$, promised to satisfy one of
\[
                 h_G\le\phi,
                 \qquad h_G\ge c\sqrt\phi;
\]
the task is to decide which alternative holds. Restrict to
$0<\phi<\min\{c^2,c^{-2},1\}$ so the alternatives are separated and
the second threshold is below one. The parameter $\phi$ is part of
the input, not a universal fixed accuracy.

NP-hardness here means a polynomial-time reduction from an NP-complete
decision problem whose outputs satisfy the promise. The question seeks
unconditional hardness at the square-root gap; hardness conditional
on another unproved conjecture is not that conclusion. The normalization
by volume is essential, since replacing it by $|S|$ on irregular graphs
defines a different expansion parameter. Disconnected graphs have
$h_G=0$ and cause no ambiguity.
\subsection{Short English statement}
Is it NP-hard to distinguish very small graph conductance from conductance at least a constant times its square root?
\subsection{Sources}
\begin{itemize}
\item[S1] ULAM.AI and the UnsolvedMath Contributors, supplied problems.json, record 2800301 (AMR-027-0301), AMR Open Problem Lists. Catalogue identity and source transcription; CC BY 4.0.
\url{https://huggingface.co/datasets/ulamai/UnsolvedMath}.
\item[S2] Bandeira - 42 Open Problems in Data Science (2016), Open Problem 3.1. Original source indexed by AMR-027-0301.
\url{https://afonsobandeira.wordpress.com/2015/10/05/10l42pcheeger/}.
\item[S3] A. S. Bandeira, Ten Lectures and Forty-Two Open Problems in the Mathematics of Data Science, corresponding numbered problem and surrounding definitions.
\url{https://people.math.ethz.ch/~abandeira/TenLecturesFortyTwoProblems.pdf}.
\end{itemize}
\end{document}
